\section{Types modernes et Pattern Matching}

Depuis C\# 8.0, le langage a introduit de nombreuses fonctionnalités visant à réduire la verbosité du code, à améliorer la sécurité à la compilation (notamment concernant les références nulles) et à encourager l'immuabilité des données.

\subsection{Types immuables : \texttt{record}, \texttt{init} et l'expression \texttt{with}}

En ingénierie logicielle, l'immuabilité garantit qu'une fois instancié, l'état d'un objet ne peut plus être modifié. Cela facilite le raisonnement sur le code, particulièrement dans des environnements concurrents. C\# simplifie l'immuabilité avec les accesseurs \texttt{init} et le mot-clé \texttt{record}.

\begin{lstlisting}[language={[Sharp]C}]
// L'accesseur 'init' permet l'affectation uniquement lors de l'instanciation
public class Person
{
    public string FirstName { get; init; }
    public string LastName { get; init; }
}

// Un 'record' est un type de reference concu pour encapsuler des donnees immuables.
// Il fournit automatiquement l'egalite par valeur, et non par reference.
public record Point(int X, int Y, int Z);

public void DemonstrateRecords()
{
    var p1 = new Point(1, 2, 3);
    var p2 = new Point(1, 2, 3);
    
    // Vrai, car les records sont compares par la valeur de leurs proprietes
    bool areEqual = (p1 == p2); 
    
    // L'expression 'with' permet de cloner un record en modifiant certaines proprietes
    var p3 = p1 with { Z = 99 };
}
\end{lstlisting}
\textit{Explication : Les records génèrent automatiquement des méthodes utiles comme \texttt{Equals}, \texttt{GetHashCode} et \texttt{ToString}. L'expression \texttt{with} effectue une copie non destructive de l'objet, en instanciant un nouveau record avec les propriétés altérées.}

\subsection{Nullable Reference Types (\texttt{string?})}

Historiquement en C\#, tous les types référence pouvaient contenir \texttt{null}, conduisant à la tristement célèbre \texttt{NullReferenceException} à l'exécution. En activant les \textit{Nullable Reference Types} (désormais actifs par défaut dans les projets modernes), le compilateur modifie cette règle : un type référence normal (\texttt{string}) ne peut plus logiquement être nul. Si le type doit supporter la nullité, il doit être déclaré explicitement avec un point d'interrogation (\texttt{string?}).

\begin{lstlisting}[language={[Sharp]C}]
// Generalement active au niveau du fichier .csproj via <Nullable>enable</Nullable>

public class UserRepository
{
    // Ne peut pas etre null. Le compilateur avertira si on tente d'y assigner null.
    public string RequiredName { get; set; } = string.Empty;

    // Peut explicitement contenir null.
    public string? OptionalMiddleName { get; set; }
    
    public void ProcessUser()
    {
        // Avertissement du compilateur : Dereferencement possible d'une reference null.
        int length = OptionalMiddleName.Length; 
        
        // Correct : Verification prealable
        if (OptionalMiddleName != null)
        {
            int safeLength = OptionalMiddleName.Length; // Plus d'avertissement
        }
    }
}
\end{lstlisting}
\textit{Explication : Cette fonctionnalité décale la gestion des erreurs de l'exécution vers la compilation, forçant le développeur à gérer la nullité explicitement via des vérifications de flux (Flow Analysis de l'analyseur statique).}

\subsection{Pattern Matching avancé (Propriétés et Tuples)}

Le Pattern Matching permet d'évaluer la forme, le type et les propriétés d'un objet de manière extrêmement concise, remplaçant avantageusement de longues chaînes de \texttt{if...else}.

\begin{lstlisting}[language={[Sharp]C}]
public record Address(string Country, string City);
public record User(string Name, Address UserAddress);

public decimal CalculateShippingCost(User user) => user switch
{
    // Property Pattern : Verification profonde des proprietes
    { UserAddress: { Country: "France", City: "Paris" } } => 5.0m,
    { UserAddress: { Country: "France" } } => 8.0m,
    
    // Declaration de variable lors du match
    { UserAddress: { Country: var country } } when country == "Belgique" || country == "Suisse" => 15.0m,
    
    // Discard pattern (attrape-tout)
    _ => 25.0m
};

// Tuple Pattern
public string GetDirection(int x, int y) => (x, y) switch
{
    (0, 0) => "Origine",
    (var a, var b) when a == b => "Diagonale",
    (_, > 0) => "Nord",
    _ => "Autre direction"
};
\end{lstlisting}
\textit{Explication : L'expression \texttt{switch} moderne retourne une valeur. Les motifs (\textit{patterns}) sont évalués de haut en bas. Le tiret bas \texttt{\_} (discard) est utilisé pour capturer tout ce qui n'a pas été intercepté.}

\subsection{La déclaration \texttt{using} moderne}

La libération des ressources non gérées (fichiers, connexions réseau) nécessite l'appel à \texttt{Dispose()}, traditionnellement géré par un bloc \texttt{using \{ ... \}}. Depuis C\# 8.0, la déclaration \texttt{using} s'affranchit des accolades.

\begin{lstlisting}[language={[Sharp]C}]
public void ReadConfigurationFile(string path)
{
    // La ressource sera liberee automatiquement a la fin de la methode.
    // Plus besoin d'accolades, ce qui evite une indentation excessive.
    using var reader = new StreamReader(path);
    
    string content = reader.ReadToEnd();
    Console.WriteLine(content);
} // <- reader.Dispose() est appele implicitement ici
\end{lstlisting}
\textit{Explication : La portée (scope) de la ressource est liée à la portée du bloc englobant (ici, la fin de la méthode). C'est syntaxiquement plus propre et réduit le niveau d'indentation du code.}

\begin{tcolorbox}[colback=red!5!white,colframe=red!75!black,title=Pièges courants \& Erreurs de compilation]
\begin{itemize}
    \item \textbf{Mauvaise utilisation de \texttt{with}} : L'expression \texttt{with} sur un \texttt{record} retourne une \textbf{nouvelle instance}. Elle ne modifie pas l'objet d'origine. Écrire \texttt{user with \{ Name = "Bob" \};} sans assigner le résultat à une variable ne sert à rien.
    \item \textbf{Ignorer les avertissements de Nullability} : Bien qu'ils ne bloquent pas toujours techniquement la compilation, ignorer les avertissements des \textit{Nullable Reference Types} détruit l'intérêt sécuritaire de la fonctionnalité.
    \item \textbf{Ordre du Pattern Matching} : L'expression \texttt{switch} évalue les cas dans l'ordre de déclaration. Placer un cas trop générique (ou le discard \texttt{\_}) au-dessus de cas spécifiques provoquera une erreur de compilation (\textit{Unreachable code}).
\end{itemize}
\end{tcolorbox}

\subsection*{Synthèse : À retenir}
\begin{itemize}
    \item Utilisez les \texttt{record} pour les objets de transfert de données (DTO) où l'égalité par valeur et l'immuabilité sont souhaitées.
    \item Utilisez l'accesseur \texttt{init} pour garantir que les propriétés ne sont mutées que lors de l'instanciation de l'objet, sécurisant son intégrité.
    \item Les expressions \texttt{switch} associées aux \textit{Property/Tuple Patterns} réduisent significativement la complexité cyclomatique du code logique.
    \item La nouvelle syntaxe \texttt{using var} sans accolades rend le code de manipulation de flux plus plat et plus lisible.
\end{itemize}
